\documentclass[11pt,a4paper]{article}
\usepackage[a4paper,margin=25mm]{geometry}
\usepackage[T1]{fontenc}
\usepackage{lmodern}
\usepackage{enumitem}
\usepackage{hyperref}

\title{GNU Backgammon by clavierhaus.at\\V0.8.11 Architecture Note}
\date{}

\begin{document}
\maketitle

\section*{Milestone}

V0.8.11 is a structural milestone.  The Android front-end now exposes a broader GNU Backgammon lifecycle and command surface through JNI, Kotlin \texttt{Engine.kt}, and \texttt{GameViewModel} wrappers.

No visible UI controls were added.  The purpose of this version is groundwork: future Android controls can now be connected to GNUbg-native commands instead of recreating game semantics in Kotlin.

\section*{Architectural Rule}

GNUbg remains authoritative for game, match, cube, session, file, and lifecycle state.  Android remains responsible for presentation, navigation, interaction, and configuration.

\section*{Exposed Command Groundwork}

The bridge groundwork covers commands such as:

\begin{itemize}[noitemsep]
  \item new game
  \item new match
  \item new session
  \item end game
  \item resign
  \item next
  \item accept, reject, decline, agree
  \item redouble
  \item load/save game
  \item load/save match
  \item load/save position
\end{itemize}

The existing SGF load/save functions remain available.  The GNUbg \texttt{quit} command is intentionally not bound at this stage because it is not part of the current Android native link and has unclear Android shell semantics.

\section*{UI Consequence}

Settings screens remain configuration-only.

Live board interactions such as roll, double, move, undo, confirm, dice handling, and destination-stack moves stay on the board surface.

The reserved lower-left area of the Play screen is now identified as a future location for true GNUbg-backed lifecycle controls, not duplicated board actions.

\section*{Status}

Android debug build and native bridge build both succeed at this milestone.

\end{document}
